\documentclass[11pt]{article}
\usepackage{graphicx} % Required for inserting images

% Preámbulo
\title{Programación II: Laboratorio 3}
\author{Diego Hernández}
\date{9 de octubre de 2025}
\setlength{\parindent}{4em}  % o el valor que prefieras

% Portada
\begin{document}
\maketitle
\newpage

% Introducción

\section{Introducción}
El siguiente informe resume contenidos introductorios a la Programación Orientada a Objetos (POO), con énfasis en el lenguaje Python.\\

\noindent La Programación Orientada a Objetos (POO) «es un paradigma de programación en el que los conceptos del mundo real relevantes para nuestro problema se modelan a través de clases y objetos, y en el que nuestro programa consiste en una serie de interacciones entre estos objetos».\\

% Desarrollo: Sección Herencia y Polimorfismo

\section{Herencia y Polimorfismo}
\subsection{Herencia}

\noindent En la POO, la \textbf{herencia} es el proceso mediante el cual una subclase hereda atributos y métodos de una súperclase o clase padre.

\subsubsection{Herencia múltiple y MRO}
Algo particular que ocurre con la herencia es que puede ocurrir la herencia múltiple, es decir, que una subclase herede métodos y atributos de más de una súperclase, ya sea porque hereda de dos súperclases de igual jerarquía o porque hereda de una súperclase que a su vez hereda de una súperclase aún mayor en jerarquía.\\

\noindent En dado caso, si se invoca algún método de una clase y esta no lo tiene en sí, sino que es parte de la herencia de sus súperclases, el orden en que debe recorrer las súperclases está definido por el \textbf{MRO} (\textit{Method Resolution Order}), el cual dictamina qué orden debe recorrer las clases padre para buscar los métodos solicitados.\\

\noindent En el caso de Python, este lenguaje a partir de su versión 2.3 usa un MRO basado en el algoritmo \textit{C3 linearization}, el cual básicamente hace que se revisen las clases inmediatamente superiores a la subclase a la que se le pide el método, y luego seguir con las superiores de un nivel mayor. Esto difiere del MRO que usaba antiguamente algoritmo \textit{DFS}, el cual hacía que la búsqueda siempre terminara en la súperclase de mayor nivel cada vez que pudiese, lo cual generaba repeticiones innecesarias.


\subsection{Polimorfismo}
\noindent Relacionado al concepto de herencia, una vez heredados los métodos y atributos en la subclase, estos pueden ser sobreescritos con características que al programador le parezca pertinente redefinir para esa subclase en particular, y estos métodos y atributos podrán conservar el mismo nombre. A esto se le llama \textbf{polimorfismo}, ya que el método o atributo es el mismo, pero tiene un funcionamiento interno (o forma) distinto.\\

\noindent Sin embargo, esto no es exclusivo de las relaciones de herencia: en el concepto de programación conocido como \textit{Duck Typing}, los métodos y atributos pueden llamarse igual en clases distintas y sin relación jerárquica, y tener formas distintas, y esto aún sería polimorfismo.\\

\noindent En resumen, el polimorfismo puede darse en todo tipo de relaciones entre clases (herencia o clases no asociadas), mientras el método o atributo se llame igual y tenga un funcionamiento distinto.\\

% Desarrollo: Sección Clases abstractas e interfaces

\section{Abstracción: Clases abstractas e Interfaces}
La \textbf{abstracción} es un término de programación y POO que hace la «ocultación de la complejidad intrínseca de una aplicación al exterior, centrándose sólo en como puede ser usada, lo que se conoce como interfaz».\\

\noindent En Python, hay formas de materializar o escribir esta abstracción a la hora de desarrollar código orientado a objetos, diseñando clases generales de las que se desprendan clases (en este sentido, subclases) con características similares entre sí, fijando atributos y métodos que deben estar sí o sí implementadas en estas subclases, pero no se define su código. Podría decirse que son «planos de construcción» de clases.\\

\noindent Hay dos modos principales en Python de llevar a cabo esta abstracción con clases generales: las clases abstractas y las interfaces.

\subsection{Clases abstractas}
Las clases abstractas se caracterizan porque permiten métodos concretos con código y se transfiere a las subclases; es decir, no es completamente abstracta.\\

\noindent En Python, usan el módulo \textit{abc} para obligar a implementar en las subclases los métodos que se declaran como métodos abstractos.

\subsection{Interfaces}

Las interfaces se diferencian de las clases abstractas en el sentido en que estas sí son completamente no implementables.\\

\noindent En Python, se usan los \textit{Protocols} del módulo \textit{typing} para verificar que un objeto cumple con el diseño de la interfaz.

\begin{itemize}
    \item En primer lugar, se diseña una interfaz y se le asigna un verificador.
    \item En segundo lugar, se crea una clase y se instancia un objeto de esa clase.
    \item Finalmente, se valida con una función si el objeto cumple con el diseño de la interfaz. Es decir, se valida el objeto en vez de clase, a diferencia de lo que sucedía con las clases abstractas.
\end{itemize}


\end{document}
